% R566: fuente francesa intro-fr.tex:24--68, pp. 5--7.
No corresponde a esta introducción dar una descripción más o menos sumaria del punto de vista de los «esquemas» en geometría algebraica, ni la larga lista de razones que han hecho necesaria su adopción y, en particular, la aceptación sistemática de elementos nilpotentes en los anillos locales de las «variedades» que consideramos (lo que necesariamente relega a un segundo plano la noción de aplicación racional en favor de la de aplicación regular o «morfismo»).
Este tratado pretende precisamente desarrollar de manera sistemática el lenguaje de los «esquemas» y demostrará, esperamos, su necesidad.
Aunque sería fácil hacerlo,

\oldpage[I]{6}
tampoco intentaremos dar aquí una introducción «intuitiva» a las nociones desarrolladas en el capítulo primero.
El lector que desee tener una visión preliminar de las materias de este tratado podrá acudir a la conferencia que A.~Grothendieck pronunció en el Congreso Internacional de Matemáticos de Edimburgo en 1958~\cite{I-7} y a la exposición~\cite{I-8} del mismo autor.
El trabajo~\cite{I-14} (citado como~(FAC)) de J.-P.~Serre también puede considerarse una exposición intermedia entre el punto de vista clásico y el de los esquemas en geometría algebraica; por ello, su lectura puede constituir una excelente preparación para la de nuestros \emph{Éléments}.

\bigskip
\begin{center}*\quad*\quad*\end{center}
\bigskip

A título informativo, damos a continuación el plan general previsto para este tratado, sujeto por lo demás a modificaciones posteriores, sobre todo en lo que respecta a los últimos capítulos:

\medskip
\begin{tabular}{r p{0.78\linewidth}}
  Capítulo primero. & --- El lenguaje de los esquemas.\\
  --- \textbf{II}. & --- Estudio global elemental de algunas clases de morfismos.\\
  --- \textbf{III}. & --- Cohomología de los haces algebraicos coherentes. Aplicaciones.\\
  --- \textbf{IV}. & --- Estudio local de los morfismos.\\
  --- \textbf{V}. & --- Procedimientos elementales de construcción de esquemas.\\
  --- \textbf{VI}. & --- Técnica de descenso. Método general de construcción de esquemas.\\
  --- \textbf{VII}. & --- Esquemas en grupos, espacios fibrados principales.\\
  --- \textbf{VIII}. & --- Estudio diferencial de los espacios fibrados.\\
  --- \textbf{IX}. & --- El grupo fundamental.\\
  --- \textbf{X}. & --- Residuos y dualidad.\\
  --- \textbf{XI}. & --- Teorías de intersección, clases de Chern, teorema de Riemann-Roch.\\
  --- \textbf{XII}. & --- Esquemas abelianos y esquemas de Picard.\\
  --- \textbf{XIII}. & --- Cohomología de Weil.\\
\end{tabular}
\medskip

En principio, todos los capítulos se consideran abiertos y siempre podrán añadírseles posteriormente párrafos suplementarios; estos aparecerán en fascículos separados, para reducir los inconvenientes del modo de publicación adoptado.
Cuando uno de esos párrafos esté previsto o en preparación al publicarse un capítulo, se mencionará en el sumario de dicho capítulo, aunque ciertas prioridades hagan que su publicación efectiva deba producirse bastante después.
Para comodidad del lector, ofrecemos en un «Capítulo~0» diversos complementos de álgebra conmutativa, álgebra homológica y teoría de haces, utilizados a lo largo de los capítulos de este tratado; son más o menos conocidos, pero no ha sido posible dar referencias prácticas para ellos.
Se recomienda al lector acudir al capítulo~0 solo durante la lectura del tratado propiamente dicho y en la medida en que los resultados a los que remitimos no le sean

\oldpage[I]{7}
suficientemente familiares.
Pensamos, además, que de esta manera la lectura del tratado podrá ser, para el principiante, un buen método de familiarizarse con el álgebra conmutativa y el álgebra homológica, cuyo estudio, cuando no va acompañado de aplicaciones tangibles, resulta tedioso e incluso deprimente para bastantes personas.

\bigskip
\begin{center}*\quad*\quad*\end{center}
\bigskip
